% Pista pro-25i-q3 · Probabilitat
% Procedència: PAU Catalunya 2025 · Convocatòria de juny (incidències) · Sèrie 4 · Exercici 3
\documentclass[11pt,a4paper]{article}
\usepackage{pista}
\pistainfo{Catalunya 2025}{Convocatòria de juny (incidències)}{Sèrie 4}

\begin{document}

\section*{\textcolor{blauPista}{Exercici 3}}

\noindent Un usuari d'Internet ha estimat que el $25\,\%$ dels correus electrònics que rep són correu brossa, mentre que la resta no ho són. Per a facilitar la classificació del correu, s'ha instal·lat un filtre que envia a la carpeta de correu brossa el $95\,\%$ dels missatges que efectivament ho són. Malauradament, aquest filtre deixa a la safata d'entrada només el $90\,\%$ dels missatges bons (i la resta els envia a la carpeta de correu brossa).

\subsection*{\textit{a)} \normalfont Quina és la probabilitat que un missatge sigui enviat pel filtre a la carpeta de correu brossa? \hspace*{\fill}\textnormal{[0,75~punts]}}

\pista{Fes un diagrama d'arbre. A la primera branca hi ha si el missatge és brossa o no; a la segona, si el filtre l'envia a la carpeta de brossa o no. Fixa't que el $90\,\%$ és dels missatges bons que es queden a la safata d'entrada. Un $90\,\%$ es tradueix en $p = 0{,}9$. Per tant, el $10\,\%$ dels missatges bons van a la carpeta de brossa.

\medskip
Hi ha dues maneres d'arribar a la carpeta de brossa: el missatge és brossa i el filtre l'hi envia, o el missatge és bo però el filtre s'equivoca i l'hi envia. Multiplica les probabilitats al llarg de cada camí i suma els dos resultats. Això és el teorema de la probabilitat total.}

\subsection*{\textit{b)} \normalfont Un dia, aquest usuari obre la carpeta de correu brossa. Quin percentatge de missatges que no són correu brossa hi trobarà? \hspace*{\fill}\textnormal{[0,75~punts]}}

\pista{Torna a l'arbre de l'apartat \textit{a)}. Saps que el missatge és a la carpeta de brossa: queda't només amb els camins que acaben en «carpeta de brossa». Et demanen quina part d'aquesta probabilitat ve del camí en què el missatge és bo:
\[
\frac{P(\text{és bo i va a la carpeta de brossa})}{P(\text{va a la carpeta de brossa})}
\]
El denominador és el resultat de l'apartat \textit{a)}. Dona la resposta en percentatge, com demana l'enunciat. Això és el teorema de Bayes.}

\vspace{0.6em}
\noindent En un servidor de correu electrònic determinat, la probabilitat de rebre com a mínim dos correus en menys de $t$ minuts ve donada per la funció
\[
F(t) = \int_{0}^{t} A\,e^{-0{,}5x}\,dx,
\]
on $A$ és una constant real.

\subsection*{\textit{c)} \normalfont Sabent que la probabilitat de rebre com a mínim dos correus en menys de dos minuts és $1 - e^{-1}$, trobeu el valor de $A$. \hspace*{\fill}\textnormal{[1~punt]}}

\pista{En primer lloc, calcula la integral definida, aplicant la regla de Barrow. Recorda aquesta primitiva:
\[
\int e^{kx}\,dx = \frac{1}{k}\,e^{kx} + C
\]
Aquí, $k = -0{,}5$. Obtindràs $F(t)$ en funció de $t$ i d'$A$.

\medskip
A continuació, imposa que $F(2) = 1 - e^{-1}$: obtindràs una equació amb la variable $A$. No és tan difícil com sembla, perquè el factor $1 - e^{-1}$ apareix als dos costats.}

\end{document}
